% ============================================================
%  Chapter 15 — Coordinate and Matrix Representation of Linear Transformations
% ============================================================
\section{Coordinate and Matrix Representation of Linear Transformations}

Every vector in a finite-dimensional space can be uniquely expressed as a linear
combination of a chosen ordered basis.  Recording these coefficients gives the
\emph{coordinate vector}, and applying this idea to a linear transformation
yields the \emph{matrix representation} — the central bridge between abstract
linear algebra and matrix computation.

% ---- Coordinate vector --------------------------------------
\begin{defbox}[Definition --- Coordinate Vector]
\label{def:coordinate-vector}
Let $V$ be a finite-dimensional vector space over $\F$ with ordered basis
$\mathcal{B}=\{v_1,\ldots,v_n\}$.  For every $v\in V$ there exist \emph{unique}
scalars $\alpha_1,\ldots,\alpha_n\in\F$ such that
\[
    v = \alpha_1 v_1 + \cdots + \alpha_n v_n.
\]
The \textbf{coordinate vector of $v$ relative to $\mathcal{B}$} is
\[
    [v]_{\mathcal{B}} =
    \begin{pmatrix}\alpha_1\\\vdots\\\alpha_n\end{pmatrix} \in \F^n.
\]
The map $[\,\cdot\,]_{\mathcal{B}}:V\to\F^n,\; v\mapsto[v]_{\mathcal{B}}$
is called the \textbf{coordinate map} relative to $\mathcal{B}$.
\end{defbox}

% ---- Coordinate map is an isomorphism -----------------------
\begin{propbox}[Proposition --- Coordinate Map is a Linear Isomorphism]
\label{prop:coord-iso}
The coordinate map $[\,\cdot\,]_{\mathcal{B}}:V\to\F^n$ is a linear
transformation and an isomorphism.  In particular:
\[
    [v+w]_{\mathcal{B}} = [v]_{\mathcal{B}} + [w]_{\mathcal{B}}, \qquad
    [cv]_{\mathcal{B}} = c\,[v]_{\mathcal{B}} \quad\text{for all }c\in\F.
\]
\end{propbox}

% ---- Examples of coordinate vectors -------------------------
\begin{exbox}[Example --- Coordinate Vectors in Standard Spaces]
\begin{enumerate}
    \item $V=\R^n$, standard basis $\varepsilon=\{e_1,\ldots,e_n\}$:
          $[x]_\varepsilon=(x_1,\ldots,x_n)^T$.
    \item $V=\F_n[x]$, basis $\varepsilon=\{1,x,\ldots,x^n\}$:
          $[a_0+a_1 x+\cdots+a_n x^n]_\varepsilon=(a_0,a_1,\ldots,a_n)^T$.
    \item $V=\R^3$, basis $\mathcal{B}=\{(1,1,1),(1,1,0),(1,0,0)\}$:
          $[(x,y,z)]_{\mathcal{B}}=(z,\,y-z,\,x-y)^T$.
\end{enumerate}
\end{exbox}

% ---- Matrix representation diagram --------------------------
\begin{flowbox}[Key Formula --- Matrix Representation Diagram]
\label{flow:matrix-rep}
Let $T:V\to W$ be linear, $\mathcal{B}=\{v_1,\ldots,v_n\}$ an ordered basis for
$V$, and $\mathcal{B}'=\{w_1,\ldots,w_m\}$ an ordered basis for $W$.
\[
\begin{array}{ccc}
V & \xrightarrow{[\,\cdot\,]_{\mathcal{B}}} & \F^n \\[6pt]
{\scriptstyle T}\Big\downarrow & & \Big\downarrow{\scriptstyle [T]^{\mathcal{B}}_{\mathcal{B}'}} \\[6pt]
W & \xrightarrow{[\,\cdot\,]_{\mathcal{B}'}} & \F^m
\end{array}
\]
The diagram commutes:
\[
    \boxed{[T(v)]_{\mathcal{B}'} = [T]^{\mathcal{B}}_{\mathcal{B}'}\,[v]_{\mathcal{B}}.}
\]
The matrix $[T]^{\mathcal{B}}_{\mathcal{B}'}$ is constructed column by column:
\[
    [T]^{\mathcal{B}}_{\mathcal{B}'} =
    \bigl[[T(v_1)]_{\mathcal{B}'}\;\big|\;[T(v_2)]_{\mathcal{B}'}\;\big|\;\cdots
          \;\big|\;[T(v_n)]_{\mathcal{B}'}\bigr].
\]
\end{flowbox}

% ---- Definition: matrix of T relative to bases --------------
\begin{defbox}[Definition --- Matrix of a Linear Transformation]
\label{def:matrix-of-lt}
Let $T:V\to W$ with ordered bases $\mathcal{B}_1=\{v_1,\ldots,v_n\}$ for $V$
and $\mathcal{B}_2=\{w_1,\ldots,w_m\}$ for $W$.  The \textbf{matrix of $T$
relative to $\mathcal{B}_1$ and $\mathcal{B}_2$} is the $m\times n$ matrix
\[
    A = [T]^{\mathcal{B}_1}_{\mathcal{B}_2}
      = \bigl[[T(v_1)]_{\mathcal{B}_2}\;\big|\;\cdots\;\big|\;
              [T(v_n)]_{\mathcal{B}_2}\bigr].
\]
When $V=W$ and $\mathcal{B}_1=\mathcal{B}_2=\mathcal{B}$, we write $[T]_\mathcal{B}$.
\end{defbox}

% ---- Unique matrix correspondence ---------------------------
\begin{thmbox}[Theorem --- Each Linear Transformation Corresponds to a Unique Matrix]
\label{thm:lt-unique-matrix}
Let $V$ be $n$-dimensional with ordered basis $\mathcal{B}$, and $W$ be
$m$-dimensional with ordered basis $\mathcal{B}'$.  For each linear
transformation $T:V\to W$ there exists a \emph{unique} matrix
$A\in\Mmn{m}{n}$ such that
\[
    [T(v)]_{\mathcal{B}'} = A\,[v]_{\mathcal{B}} \quad \text{for all }v\in V.
\]
Moreover, the assignment $T\mapsto A=[T]^{\mathcal{B}}_{\mathcal{B}'}$ is a
one-to-one correspondence between $\Hom_\F(V,W)$ and $\Mmn{m}{n}$.
\end{thmbox}

% ---- Kernel/image via matrix representation -----------------
\begin{corbox}[Corollary --- Kernel and Image via Matrix Representation]
\label{cor:ker-img-matrix}
Let $T:V\to W$ with bases $\mathcal{B},\mathcal{B}'$ and matrix
$A=[T]^{\mathcal{B}}_{\mathcal{B}'}$.
\begin{enumerate}[label=(\roman*)]
    \item $\Ker T = \{v\in V \mid [v]_{\mathcal{B}}\in\Ker A\}$.
          A set $\{u_1,\ldots,u_r\}$ is a basis of $\Ker T$ if and only if
          $\{[u_1]_{\mathcal{B}},\ldots,[u_r]_{\mathcal{B}}\}$ is a basis
          of $\Ker A$.
    \item $\Img T = \{w\in W \mid [w]_{\mathcal{B}'}\in\operatorname{Col}(A)\}$.
          A set $\{u_1,\ldots,u_s\}$ is a basis of $\Img T$ if and only if
          $\{[u_1]_{\mathcal{B}'},\ldots,[u_s]_{\mathcal{B}'}\}$ is a basis
          of $\operatorname{Col}(A)$.
\end{enumerate}
\end{corbox}

% ---- Addition and scalar mult of l.t. -----------------------
\begin{tipbox}[Remark --- Addition and Scalar Multiplication of Linear Transformations]
\label{rem:lt-add-scalar}
For linear transformations $T,T':V\to W$ and scalar $\alpha\in\F$:
\[
    [\alpha T + T']^{\mathcal{B}}_{\mathcal{B}'}
    = \alpha\,[T]^{\mathcal{B}}_{\mathcal{B}'} + [T']^{\mathcal{B}}_{\mathcal{B}'}.
\]
Hence the matrix representation is compatible with the vector space structure
on $\Hom_\F(V,W)$.
\end{tipbox}

% ---- F_{BB'} is an isomorphism ------------------------------
\begin{thmbox}[Theorem --- Matrix Representation is an Isomorphism of Vector Spaces]
\label{thm:matrix-rep-iso}
The map
\[
    \Phi_{\mathcal{B}\mathcal{B}'}: \Hom_\F(V,W) \longrightarrow \Mmn{m}{n},
    \quad T \longmapsto [T]^{\mathcal{B}}_{\mathcal{B}'},
\]
is an isomorphism of vector spaces over $\F$.
\end{thmbox}

\begin{corbox}[Corollary --- Dimension of $\Hom_\F(V,W)$]
\label{cor:dim-hom}
If $\dime_\F V = n$ and $\dime_\F W = m$, then
\[
    \dime_\F\bigl(\Hom_\F(V,W)\bigr) = mn.
\]
\end{corbox}

% ---- Composition and matrix product -------------------------
\begin{thmbox}[Theorem --- Composition Corresponds to Matrix Product]
\label{thm:comp-matrix-product}
Let $T_1:V\to W$ and $T_2:W\to Z$ be linear transformations with ordered
bases $\mathcal{B}_1,\mathcal{B}_2,\mathcal{B}_3$ for $V,W,Z$ respectively.
Then
\[
    [T_2\circ T_1]^{\mathcal{B}_1}_{\mathcal{B}_3}
    = [T_2]^{\mathcal{B}_2}_{\mathcal{B}_3}\,[T_1]^{\mathcal{B}_1}_{\mathcal{B}_2}.
\]
\end{thmbox}

% ---- Invertible l.t. <-> invertible matrix ------------------
\begin{corbox}[Corollary --- Invertible Linear Map Corresponds to Invertible Matrix]
\label{cor:invertible-lt-matrix}
Let $T:V\to W$ be invertible with ordered bases $\mathcal{B}_1$ for $V$ and
$\mathcal{B}_2$ for $W$.  Then $[T]^{\mathcal{B}_1}_{\mathcal{B}_2}$ is
invertible and
\[
    \bigl[T^{-1}\bigr]^{\mathcal{B}_2}_{\mathcal{B}_1}
    = \bigl([T]^{\mathcal{B}_1}_{\mathcal{B}_2}\bigr)^{-1}.
\]
\end{corbox}

% ---- Change of basis ----------------------------------------
\begin{defbox}[Definition --- Change-of-Basis Matrix]
\label{def:cob-matrix}
Let $V$ be $n$-dimensional with two ordered bases
$\mathcal{B}_1=\{v_1,\ldots,v_n\}$ and $\mathcal{B}_2=\{w_1,\ldots,w_n\}$.
For each $j$ write $w_j = \sum_i P_{ij}\,v_i$.  The \textbf{change-of-basis
matrix from $\mathcal{B}_2$ to $\mathcal{B}_1$} is
\[
    C(\mathcal{B}_2,\mathcal{B}_1)
    = [\mathrm{id}]^{\mathcal{B}_2}_{\mathcal{B}_1}
    = \bigl[[w_1]_{\mathcal{B}_1}\;\big|\;[w_2]_{\mathcal{B}_1}\;\big|\;\cdots
             \;\big|\;[w_n]_{\mathcal{B}_1}\bigr].
\]
It satisfies
\[
    [v]_{\mathcal{B}_1} = C(\mathcal{B}_2,\mathcal{B}_1)\,[v]_{\mathcal{B}_2}
    \quad\text{for all }v\in V,
\]
and is invertible with $C(\mathcal{B}_2,\mathcal{B}_1)^{-1}=C(\mathcal{B}_1,\mathcal{B}_2)$.
\end{defbox}

% ---- Theorems about change of basis -------------------------
\begin{thmbox}[Theorem --- Unique Change-of-Basis Matrix]
\label{thm:cob-unique}
There exists a unique invertible matrix $A\in\Mn{n}$ such that
\[
    [v]_{\mathcal{B}_1} = A\,[v]_{\mathcal{B}_2}
    \quad\text{and}\quad
    [v]_{\mathcal{B}_2} = A^{-1}[v]_{\mathcal{B}_1}
    \quad\text{for all }v\in V.
\]
This matrix is $A = C(\mathcal{B}_2,\mathcal{B}_1)$.
\end{thmbox}

\begin{thmbox}[Theorem --- Unique Basis for a Given Change-of-Basis Matrix]
\label{thm:cob-unique-basis}
Given an invertible matrix $A\in\Mn{n}$ and an ordered basis $\mathcal{B}_1$
of $V$, there exists a unique ordered basis $\mathcal{B}_2$ of $V$ such that
$[v]_{\mathcal{B}_1} = A\,[v]_{\mathcal{B}_2}$ for all $v\in V$.
\end{thmbox}

\begin{propbox}[Proposition --- Composition of Change-of-Basis Matrices]
\label{prop:cob-composition}
For three ordered bases $\mathcal{B}_1,\mathcal{B}_2,\mathcal{B}_3$ of $V$:
\[
    C(\mathcal{B}_1,\mathcal{B}_3) = C(\mathcal{B}_2,\mathcal{B}_3)\,C(\mathcal{B}_1,\mathcal{B}_2).
\]
\end{propbox}

% ---- Change of basis for l.t. representation ----------------
\begin{thmbox}[Theorem --- Change of Basis for Matrix Representation of a Linear Map]
\label{thm:cob-lt}
Let $T:V\to W$ be linear, $\mathcal{B}_1,\mathcal{B}_2$ ordered bases for $V$,
and $\mathcal{B}'_1,\mathcal{B}'_2$ ordered bases for $W$.  Then
\[
    [T]^{\mathcal{B}_2}_{\mathcal{B}'_2}
    = C(\mathcal{B}'_1,\mathcal{B}'_2)\,[T]^{\mathcal{B}_1}_{\mathcal{B}'_1}\,
      C(\mathcal{B}_2,\mathcal{B}_1).
\]
\end{thmbox}

% ---- Endomorphism change of basis ---------------------------
\begin{corbox}[Corollary --- Change of Basis for an Endomorphism]
\label{cor:cob-endo}
For an endomorphism $T:V\to V$ and two ordered bases $\mathcal{B}_1,\mathcal{B}_2$:
\[
    [T]_{\mathcal{B}_2}
    = C(\mathcal{B}_2,\mathcal{B}_1)^{-1}\,[T]_{\mathcal{B}_1}\,
      C(\mathcal{B}_2,\mathcal{B}_1).
\]
In particular, $[T]_{\mathcal{B}_1}$ and $[T]_{\mathcal{B}_2}$ are
\textbf{similar matrices}.
\end{corbox}

\begin{tipbox}[Remark --- Geometric Meaning of Similarity]
Two matrices represent the same linear transformation in different coordinate
systems if and only if they are similar.  The transformation itself does not
change; only the way we \emph{describe} it in coordinates changes.
\end{tipbox}

% ---- Similar matrices ---------------------------------------
\begin{defbox}[Definition --- Similar Matrices]
\label{def:similar}
A matrix $B\in\Mn{n}$ is \textbf{similar} to $A\in\Mn{n}$, written
$A\sim_{\mathrm{sim}} B$, if there exists an invertible matrix $C\in\Mn{n}$
such that
\[
    B = C^{-1}AC.
\]
\end{defbox}

\begin{propbox}[Proposition --- Similarity is an Equivalence Relation]
\label{prop:similarity-equivalence}
Matrix similarity is an equivalence relation on $\Mn{n}$:
\begin{enumerate}[label=(\roman*)]
    \item \emph{Reflexive:} $A\sim A$ (take $C=I$).
    \item \emph{Symmetric:} $A\sim B\Rightarrow B\sim A$.
    \item \emph{Transitive:} $A\sim B$ and $B\sim C\Rightarrow A\sim C$.
\end{enumerate}
\end{propbox}
